% from ega4-13-fr.tex
\oldpage[IV-3]{194}
composantes irréductibles de $f^{-1}(y)$, $z_i$ le point générique de $Z_i$,
et considérons les conditions suivantes~:
\begin{enumerate}[label=\emph{\alph*)}]
  \item Pour tout $x\in f^{-1}(y)$, on a la relation
  \[
  \tag{13.2.10.1}
  \label{IV.13.2.10.1-fr}
    \dim(\mathcal O_{X,x})=\dim(\mathcal O_{Y,y})
    +\dim(\mathcal O_{X,x}\otimes_{\mathcal O_{Y,y}}k(y)).
  \]
  \item Pour tout $i$, on a
  \[
  \tag{13.2.10.2}
  \label{IV.13.2.10.2-fr}
    \dim(\mathcal O_{X,z_i})=\dim(\mathcal O_{Y,y}).
  \]
  \item Pour tout $i$, il existe une composante irréductible $X_i$ de $X$
  contenant $Z_i$ et telle que
  $\dim(X_i\cap f^{-1}(\eta))=\dim(Z_i)$ (autrement dit, telle que le
  \emph{sous-préschéma fermé réduit} $X_i$ de $X$ soit équidimensionnel sur $Y$
  au point $z_i$).
\end{enumerate}
Alors $a)$ entraîne $b)$ et $b)$ entraîne $c)$~; en outre, si
$\mathcal O_{Y,y}$ est universellement caténaire, les trois conditions $a)$,
$b)$, $c)$ sont équivalentes.
\end{proposition}

L'anneau $\mathcal O_{X,z_i}\otimes_{\mathcal O_{Y,y}}k(y)$ étant de dimension
$0$ (5.1.2), $a)$ entraîne évidemment $b)$~; $b)$ entraîne $c)$ en vertu de
(13.2.7) appliqué au point $z_i$. Inversement, supposons que
$\mathcal O_{Y,y}$ soit universellement caténaire~; comme la condition $c)$
implique que $f(X_i)$ est dense dans $Y$, il résulte de $c)$ que les conditions
(i) à (iv) de (13.2.7) sont vérifiées en remplaçant $X'$ par $X_i$ et $x$ par
$z_i$, donc $c)$ implique $b)$~; enfin $b)$ implique $a)$ en vertu de
(13.2.9).

\begin{corollary}[13.2.11]
\label{IV.13.2.11-fr}
Les notations étant celles de (13.2.10), supposons que $\mathcal O_{Y,y}$ soit
universellement caténaire. Pour tout $y'\in Y$, soit $E(y')$ l'ensemble des
dimensions des composantes irréductibles de $f^{-1}(y')$ et posons
$d(y')=\dim(f^{-1}(y'))=\sup(E(y'))$. Alors, si les conditions équivalentes
$a)$, $b)$, $c)$ de (13.2.10) sont vérifiées, on a
$E(y)\subset E(\eta)$, d'où $d(y)\leq d(\eta)$.
\end{corollary}

En effet, avec les notations de (13.2.10), $X_i\cap f^{-1}(\eta)$ n'est pas
vide et est par suite une composante irréductible de $f^{-1}(\eta)$
($\mathbf 0_{\mathrm I}$, 2.1.8).

\begin{remarks}[13.2.12]
\label{IV.13.2.12-fr}
\begin{enumerate}[label=(\roman*)]
  \item Rappelons (6.1.2) que la relation (13.2.10.1) est toujours vérifiée
  lorsque le morphisme $f$ est \emph{plat} au point $x$.
  \item Avec les notations de (13.2.10), supposons que $\mathcal O_{Y,y}$ soit
  universellement caténaire et en outre que le morphisme $f$ soit \emph{propre}~;
  alors, si les conditions équivalentes $a)$, $b)$, $c)$ de (13.2.10) sont
  satisfaites, on a même $d(y)=d(\eta)$, car il résulte de (13.1.5) que l'on a
  $d(y)\geq d(\eta)$.
  \item Le morphisme de (12.2.3, (ii)) est \emph{propre} et \emph{plat} et tous
  les anneaux locaux de $Y$ sont universellement caténaires~; en outre, les deux
  composantes irréductibles $X_1$, $X_2$ de $X$ sont \emph{équidimensionnelles}
  sur $Y$ en tout point~; mais $E(y)$ a \emph{deux} éléments pour tout
  $y\neq y_0$, alors que $E(y_0)$ est réduit à un \emph{seul} élément, donc
  $E(y)$ n'est \emph{pas constant} dans $Y$.
\end{enumerate}
\end{remarks}

\subsection{Morphismes équidimensionnels : cas général}
\label{subsection:IV.13.3-fr}

\begin{proposition}[13.3.1]
\label{IV.13.3.1-fr}
Soient $Y$ un préschéma, $f:X\to Y$ un morphisme localement de type fini,
$x$ un point de $X$, $y=f(x)$. Désignons par $Y_\alpha$ les composantes
irréductibles de $Y$ contenant $y$. Alors les conditions suivantes sont
équivalentes~:
\begin{enumerate}[label=\emph{\alph*)}]
